\chapter{Phase Transitions in Language Model Output}
\label{ch15}

The temperature parameter of a language model is not merely a knob for controlling
output randomness. At certain critical values it can trigger qualitative, system-wide
changes in the structure of the output distribution --- changes that physicists call
\emph{phase transitions}. This chapter makes that claim precise, connects it to the
singularity theory developed in Chapters 1--5, and reviews the growing empirical
evidence from large-scale experiments on GPT-2, Pythia, Mistral, and Llama.

\section{Phase Transitions: The Analytic Definition}
\label{ch15:analytic-definition}

Throughout Chapters 13--14 we associated to any probability measure $\mu$ on
$\Sigma^*$ the \emph{partition function}
\[
  Z(\beta) = \sum_{w \in \Sigma^*} \mu(w)^{\,\beta}, \qquad \beta > 0,
\]
and the \emph{free energy}
\[
  f(\beta) = -\frac{1}{\beta}\log Z(\beta).
\]
The Gibbs measure at inverse temperature $\beta$ assigns to each word $w$ the
probability $\mu_\beta(w) = \mu(w)^\beta / Z(\beta)$.  Temperature in the sense
of Chapter 14 is $T = 1/\beta$.

\begin{definition}[Phase transition]
\label{ch15:def-phase-transition}
A \emph{phase transition} of the system $(\mu, \beta)$ is a point $\beta_c \in
(0,\infty)$ at which the free energy $f(\beta)$ fails to be real-analytic. The
transition is \emph{first-order} if $f'$ has a jump discontinuity at $\beta_c$;
it is \emph{second-order} (or \emph{continuous}) if $f'$ is continuous at
$\beta_c$ but $f''$ diverges there.
\end{definition}

In statistical mechanics the second derivative $-\beta^2 f''(\beta)$ is
proportional to the \emph{susceptibility}, measuring how sensitively the mean
energy responds to a change in temperature. A diverging susceptibility is
therefore the signature of a second-order transition.

\subsection{Connection to Singularities of Generating Functions}
\label{ch15:gf-connection}

Write $z = e^{-\beta}$, so $Z(\beta) = A(z)$ where $A(z) = \sum_{w} \mu(w)^\beta$
is interpreted as a generating function evaluated at $z = e^{-\beta}$.  Then
$\beta_c$ is a phase-transition point if and only if $z_c = e^{-\beta_c}$ is a
singularity of $A(z)$ on the boundary of its disk of convergence.

This is precisely the setting of Chapter 4's singularity classification:

\begin{itemize}
  \item A \emph{simple pole} of $A(z)$ at $z_c$ gives a first-order transition.
    The dominant coefficient grows as $z_c^{-n}$ with no polynomial correction,
    and the free energy has a kink.
  \item A \emph{square-root branch point}, $A(z) \sim C(1 - z/z_c)^{1/2}$,
    gives a second-order transition: $Z(\beta)$ is continuous through $\beta_c$
    but $Z''(\beta)$ diverges.  The associated coefficient asymptotics carry the
    $n^{-3/2}$ factor familiar from Chapter 5.
  \item An \emph{essential singularity} can produce more exotic behavior
    (infinitely many derivatives discontinuous) that lies outside the standard
    first/second-order classification.
\end{itemize}

\begin{proposition}
\label{ch15:2nd-order}
Suppose $A(z)$ has a unique dominant singularity at $z_c$ of the square-root
algebraic type,
\[
  A(z) = g(z) - h(z)\sqrt{1 - z/z_c},
\]
where $g, h$ are analytic and nonzero at $z_c$. Then the free energy
$f(\beta) = -\beta^{-1}\log A(e^{-\beta})$ is $C^1$ at $\beta_c = -\log z_c$
but $f''$ has an integrable singularity there, giving a second-order phase
transition.
\end{proposition}

The proof is a direct computation: one writes
$A(e^{-\beta}) = g(e^{-\beta}) - h(e^{-\beta})\sqrt{1 - e^{-(\beta-\beta_c)}}$
and expands near $\beta = \beta_c$.  The square root produces a branch cut in
$\beta$, so all higher derivatives of $f$ diverge at $\beta_c$, while $f$ and $f'$
remain finite.

The upshot: \emph{the order of the phase transition is determined by the type of
the dominant singularity of $A(z)$}.  This is the bridge between Chapter 4's
taxonomy and the thermodynamics of Chapters 13--14.

\section{Composition Schemes and Universal Transition Signatures}
\label{ch15:composition-schemes}

A natural mechanism by which a smooth generating function can develop a
singularity is the \emph{composition scheme}: take an outer function $A(u)$,
analytic in a neighborhood of $u=0$, and an inner function $B(z)$ with radius of
convergence $\rho_B$.  The composed function $F(z) = A(B(z))$ has its dominant
singularity at $z = \rho_B$, but the \emph{type} of that singularity depends on
whether $\rho_B$ is an ordinary point or a critical point of $A$.

\cite{BanderierKubaWallner2021} carried out a systematic analysis of such
composition schemes and classified the possible singularity types that arise. The
central finding is:

\begin{itemize}
  \item If $B(\rho_B) < \rho_A$ (where $\rho_A$ is the radius of convergence of
    $A$), the singularity of $F$ at $\rho_B$ is of the same type as that of $B$.
  \item If $B(\rho_B) = \rho_A$ --- the \emph{critical} or \emph{confluent} case
    --- the singularity type changes, typically from rational to square-root
    algebraic.  The coefficient asymptotics transition from exponential-only to
    exponential times $n^{-3/2}$, with an Airy-function limit law governing the
    fluctuations.
  \item More degenerate confluences (higher-order critical points) produce
    $n^{-5/3}$ and related exponents tied to higher Airy-type kernels.
\end{itemize}

In the language of this chapter, the passage from the sub-critical to the
critical composition scheme is itself a second-order phase transition.  As a
control parameter (say, the mean branching number in a context-free grammar)
crosses a threshold, the dominant singularity of $F(z)$ changes character, and
the free energy acquires a new non-analyticity.

\begin{remark}
\label{ch15:remark-universality}
The Airy-function limit law is \emph{universal} in the sense that the same law
emerges for a wide variety of combinatorial structures (binary trees, planar maps,
context-free languages) whenever the composition scheme is at criticality.  This
universality is the analytic-combinatorics analog of universality classes in
physics.
\end{remark}

The composition-scheme framework gives a \emph{catalogue} of possible phase-transition
signatures for systems whose partition function is (or is approximated by) a composed
generating function.  Chapter 18 will revisit which entries of this catalogue are
plausible candidates for real language models.

\section{A Worked Example: The Catalan Partition Function}
\label{ch15:catalan-example}

Before turning to the empirical literature, we work through the simplest non-trivial
case.

\begin{example}[Catalan partition function]
\label{ch15:ex-catalan}
Let $A(z) = \frac{1 - \sqrt{1-4z}}{2z}$, the generating function of the Catalan
numbers $C_n = \binom{2n}{n}/(n+1)$.  Its radius of convergence is $\rho = 1/4$,
with a square-root branch point there.

Define the partition function $Z(\beta) = A(e^{-\beta})$ for $\beta > 0$.  The
critical inverse temperature is
\[
  \beta_c = -\log \rho = -\log(1/4) = \log 4.
\]
For $\beta > \beta_c$ we have $e^{-\beta} < 1/4$ and $Z(\beta)$ is analytic; for
$\beta < \beta_c$ the argument exits the radius of convergence and $Z(\beta)$ is
defined by analytic continuation through the cut.

Near $\beta_c$ write $\epsilon = \beta - \beta_c$ and expand:
\[
  1 - 4e^{-\beta} = 1 - 4e^{-\beta_c - \epsilon} = 1 - e^{-\epsilon}
  \approx \epsilon - \tfrac{\epsilon^2}{2} + O(\epsilon^3).
\]
Thus
\[
  Z(\beta) \approx \frac{1 - \sqrt{\epsilon}}{2e^{-\beta_c}} + O(\epsilon)
  = 2(1 - \sqrt{\epsilon} + O(\epsilon))
\]
(where we use $e^{\beta_c} = 4$ and the leading term of the Catalan GF at the
singularity).  The free energy is
\[
  f(\beta) = -\frac{1}{\beta}\log Z(\beta) \approx -\frac{1}{\beta_c}
  \bigl(\log 2 - \tfrac{1}{2}\epsilon + O(\epsilon)\bigr).
\]
Computing: $f'(\beta)$ is continuous at $\beta_c$ (the $\sqrt{\epsilon}$ term in
$Z$ contributes a $\epsilon^{-1/2}$ term to $Z'/Z$, which enters $f'$ as
$(\log Z)' / \beta + (\log Z) / \beta^2$; the leading $\epsilon^{-1/2}$
singularity is integrable). However, $f''(\beta)$ diverges as $\epsilon \to 0$
like $\epsilon^{-1/2}$, confirming a second-order transition.

The coefficient asymptotics of $A(z)$, namely $[z^n]A(z) = C_n \sim
4^n n^{-3/2} / \sqrt{\pi}$, appear on the Gibbs side as the probability mass
that the Boltzmann sampler assigns to strings of length $n$ at the critical
temperature: it decays as $n^{-3/2}$, the signature of the algebraic singularity
noted in Chapter 5.
\end{example}

This example is the canonical prototype.  Every algebraic generating function
with a square-root dominant singularity gives the same exponents, which is why
the $n^{-3/2}$ law and its associated Airy limit are called \emph{universal}.

\section{Empirical Evidence from Language Models}
\label{ch15:empirical}

We now turn to recent experimental work.  The experiments are independent and use
different models, but they converge on a consistent picture: GPT-family language
models exhibit temperature-driven phase transitions whose signatures are
quantitatively consistent with second-order universality.

\subsection{GPT-2: Nakaishi, Nishikawa, and Hukushima (2024)}
\label{ch15:nakaishi}

\cite{NakaishiNishikawaHukushima2024} (arXiv:2406.05335) studied GPT-2 by
treating the temperature $T = 1/\beta$ as a tunable control parameter and
measuring standard finite-size-scaling observables on the output distribution.
Their methodology is borrowed from computational statistical mechanics: they
measure a quantity analogous to the magnetic susceptibility, locate its peak as a
function of $T$, and track how the peak height and location shift as the model
size (or the prompt length used as a proxy for system size) varies.

Their main findings are:

\begin{itemize}
  \item There is a temperature $T_c \approx 1$ (i.e., $\beta_c \approx 1$) at
    which the susceptibility-analog peaks sharply.
  \item The peak height scales with the system-size proxy in a manner consistent
    with a second-order transition, with fitted critical exponents.
  \item The transition separates a low-temperature \emph{ordered} phase (where
    the output concentrates on a small number of high-probability completions) from
    a high-temperature \emph{disordered} phase (where the output approaches
    uniform).
\end{itemize}

These are empirical observations with fitted exponents.  The paper does not prove
that $Z(\beta)$ is non-analytic at $\beta_c$; it demonstrates scaling behavior
that is consistent with such non-analyticity.  The distinction matters: finite
models always have analytic $Z(\beta)$ (the sum is finite), so what one observes
in practice is a rapid crossover that sharpens as model size grows, analogous to
the finite-size rounding of phase transitions in statistical mechanics.

\subsection{Pythia, Mistral, Llama: Arnold et al. (2024)}
\label{ch15:arnold}

\cite{ArnoldEtAl2024} (arXiv:2405.17088) extended the investigation to three
open-weight model families: Pythia (70M to 12B parameters), Mistral (7B), and
Llama (7B to 70B). They scanned temperature and measured several observables
including token-rank distributions, entropy, and a susceptibility proxy across
all these models.

The key findings of \cite{ArnoldEtAl2024} are:

\begin{itemize}
  \item Phase-transition signatures are \emph{robust across model families and
    parameter scales}.  The critical temperature $T_c$ varies somewhat, but all
    models show qualitatively identical behavior.
  \item The scaling of observables near $T_c$ is consistent with second-order
    criticality.
  \item There is a \emph{temperature window} around $T = 1$ where the output
    distribution is neither too concentrated nor too flat, and within this window
    the Zipf-like rank-frequency exponent is approximately constant across models.
\end{itemize}

The robustness across model families is significant: it suggests the phenomenon
is not an artifact of a particular training procedure or architecture, but a
structural property of the output distributions produced by autoregressive models
trained on natural-language corpora.

Again, this is empirical evidence.  The universality class is inferred from
exponent fits, not derived from first principles.

\subsection{Universality and the Temperature Window: Mikhaylovskiy (2025)}
\label{ch15:mikhaylovskiy}

\cite{Mikhaylovskiy2025} goes a step further and argues that the observed
transitions across models sit in a \emph{common universality class} that is
consistent with the predictions of analytic-combinatorics composition schemes.
Specifically, the fitted critical exponents are compatible with those arising from
square-root singularities of composed generating functions --- the same universality
class as the Catalan example of Section~\ref{ch15:catalan-example}.

\cite{Mikhaylovskiy2025} also identifies and analyzes the \emph{temperature window}
regime: a range of temperatures near $T_c$ in which the rank-frequency (Zipf)
exponent $\alpha$ satisfies $\alpha \approx 1$ and is approximately constant as a
function of $T$.  This is consistent with the plateau in the free energy's second
derivative that one would expect if the singularity is approached slowly from
a non-analytic direction.  Chapter 16 will examine the temperature window in
detail; here we note only that its existence is further evidence of an underlying
analytic structure.

The universality claim of \cite{Mikhaylovskiy2025} is stronger than those of
\cite{NakaishiNishikawaHukushima2024} and \cite{ArnoldEtAl2024} --- it makes a
claim about \emph{which} universality class, not merely that one exists.  It
remains to be verified whether the exponents are truly universal or whether
finite-size effects or model-family-specific features contaminate the fits.

\subsection{Additional Corroboration: Li et al. (2024)}
\label{ch15:li}

\cite{LiEtAl2024} provide additional empirical corroboration from a somewhat
different angle.  They carry out a finite-size-scaling analysis of LLM output
distributions and show that an $O(N)$ estimation procedure for the partition
function (where $N$ is vocabulary size or sequence length) reproduces the
transition signatures found by the other groups.  Their work is notable for its
computational efficiency: rather than sampling extensively, they use analytic
approximations to $Z(\beta)$ that are tractable for large models, making the
phase-transition analysis practical at scale.

\section{The Bridging Proposal}
\label{ch15:bridge}

We can now state the central proposal of Part V of this book.

\begin{remark}[The analytic-combinatorics explanation]
\label{ch15:remark-bridge}
The empirical phase transitions observed by
\cite{NakaishiNishikawaHukushima2024}, \cite{ArnoldEtAl2024}, and
\cite{Mikhaylovskiy2025} are proposed to be \emph{analytic phase transitions}
of the partition function
\[
  Z(\beta) = \sum_{w \in \Sigma^*} \mu(w)^{\,\beta},
\]
where $\mu$ is the LLM's output distribution, viewed through the Boltzmann
correspondence of Chapter 13 and the singularity classification of Chapter 4.
\end{remark}

To make this concrete, consider what happens if $\mu$ were \emph{exactly} a
PCFG distribution in the sense of Chapter 8.  Then by the Chomsky--Schützenberger
theorem (Chapter 5), the generating function $A(z) = Z(\beta)\big|_{z=e^{-\beta}}$
would be algebraic over $\mathbb{Q}(z)$.  Its dominant singularity would be a
branch point of algebraic type, typically square-root.  The partition function
$Z(\beta) = A(e^{-\beta})$ would then exhibit a second-order phase transition at
$\beta_c = -\log \rho$, with universal $n^{-3/2}$ Boltzmann weights at
criticality.  The Nakaishi--Arnold--Mikhaylovskiy scaling would be a theorem
rather than an empirical observation.

The LLM case is not provably of this form.  A trained Transformer is not a PCFG;
its output distribution has no known closed-form algebraic structure.  But the
empirical scaling is consistent with the algebraic scenario, and the robustness
across model families suggests that the phenomenon is structural rather than
incidental.

The working hypothesis, which Chapter 18 will list as Gap 2 among the ten open
problems of Part V, is:

\begin{quote}
\emph{The output distribution $\mu$ of a large autoregressive language model
trained on natural language admits an effective algebraic approximation in the
sense that $Z(\beta) = \sum_w \mu(w)^\beta$ has a dominant singularity of
square-root type at some $\beta_c$, and the resulting second-order phase
transition is what the empirical experiments are detecting.}
\end{quote}

What would it take to verify this?  One would need either (a) a theoretical
derivation of the algebraic structure of the LLM's output distribution --- which
seems currently out of reach --- or (b) a more refined experimental protocol that
distinguishes square-root from other singularity types by measuring not just the
location but the \emph{shape} of the transition in detail.  Chapter 16 takes
a step toward (b) by analyzing the temperature window, and Chapter 19 discusses
what theoretical progress on (a) might look like.

\section{Summary}
\label{ch15:summary}

The main points of this chapter:

\begin{enumerate}
  \item A phase transition of a temperature-parameterized system is a
    non-analyticity of the free energy $f(\beta)$. The order of the transition is
    determined by which derivative of $f$ first diverges.
  \item Via the substitution $z = e^{-\beta}$, the type of the dominant
    singularity of $A(z) = Z(\beta)$ determines the order of the transition.
    Square-root singularities give second-order transitions; poles give
    first-order.
  \item Composition schemes \cite{BanderierKubaWallner2021} provide a catalogue
    of universal transition signatures, with the Airy-function universality class
    (associated with $n^{-3/2}$ coefficient asymptotics) being the most
    prominent.
  \item Empirical experiments on GPT-2
    \cite{NakaishiNishikawaHukushima2024}, and on Pythia, Mistral, and Llama
    \cite{ArnoldEtAl2024}, provide evidence for temperature-driven phase
    transitions consistent with second-order criticality. \cite{Mikhaylovskiy2025}
    identifies the likely universality class, and \cite{LiEtAl2024} provide
    additional corroboration.
  \item The bridging proposal: these empirical transitions may be explained as
    analytic phase transitions of the LLM partition function, consistent with the
    PCFG/algebraic scenario but not yet proved to be so.
\end{enumerate}

\section{Notes and Further Reading}
\label{ch15:notes}

The connection between singularities of generating functions and thermodynamic
phase transitions is classical in the theory of exactly solvable models; see
\cite{FlajoletSedgewick2009}, Chapter~VII, for the combinatorics perspective.
Composition schemes and their singularity-confluence phenomena are developed
in detail in \cite{BanderierKubaWallner2021}; the Airy limit law they discuss
has a long history in map enumeration going back to Tutte.

The statistical-mechanics methodology applied to language models
--- finite-size scaling, susceptibility peaks, critical exponent fitting --- is
standard in the physics of disordered systems.  Readers unfamiliar with this
methodology should consult any graduate statistical-mechanics text before trying
to interpret the exponent fits in
\cite{NakaishiNishikawaHukushima2024} and \cite{ArnoldEtAl2024}.

The temperature window and its relation to Zipf's law in LLM output is discussed
further in Chapter 16.  The relationship between the ARM--EBM bijection
(Chapter 14) and the phase-transition picture deserves further study: briefly,
the EBM formulation makes it natural to define an order parameter for the
transition, which the ARM formulation does not provide directly.

\begin{exercise}
\label{ch15:ex1}
Let $A(z) = (1-z)^{-2}$, which has a double pole at $z=1$. Compute $Z(\beta) =
A(e^{-\beta})$ for $\beta > 0$, identify $\beta_c$, and compute $f(\beta) =
-\beta^{-1}\log Z(\beta)$ near $\beta_c$. Show that $f'$ has a jump discontinuity
at $\beta_c$, confirming a first-order transition. Compare with the Catalan
(square-root) case of Section~\ref{ch15:catalan-example}.
\end{exercise}

\begin{exercise}
\label{ch15:ex2}
The generating function of Motzkin numbers satisfies $M(z) = 1 + z M(z) + z^2
M(z)^2$, giving $M(z) = (1 - z - \sqrt{(1-z)^2 - 4z^2}) / (2z^2)$ with radius of
convergence $\rho = 1/3$. Let $Z_M(\beta) = M(e^{-\beta})$ and $\beta_c = \log 3$.
(a) Show that $Z_M$ has a square-root singularity at $\beta_c$, giving a
second-order transition. (b) Determine the coefficient asymptotics $[z^n]M(z)$
and verify that the Boltzmann weight at $\beta_c$ decays as $n^{-3/2}$. (c)
Explain why the Catalan and Motzkin systems are in the same universality class
despite having different $\rho$ and different $\beta_c$.
\end{exercise}
